
\documentclass{article}
\usepackage{colortbl}
\usepackage{makecell}
\usepackage{multirow}
\usepackage{supertabular}

\begin{document}

\newcounter{utterance}

\centering \large Interaction Transcript for game `cladder', experiment `full\_v1.5\_default', episode 3134 with Qwen3.5{-}9B.
\vspace{24pt}

{ \footnotesize  \setcounter{utterance}{1}
\setlength{\tabcolsep}{0pt}
\begin{supertabular}{c@{$\;$}|p{.15\linewidth}@{}p{.15\linewidth}p{.15\linewidth}p{.15\linewidth}p{.15\linewidth}p{.15\linewidth}}
    \# & \multicolumn{2}{c}{Player} && \multicolumn{2}{c}{Game Master} \\
    \hline

    \theutterance \stepcounter{utterance}  
    & & & \multicolumn{4}{p{0.6\linewidth}}{
        \cellcolor[rgb]{0.9,0.9,0.9}{
            \makecell[{{p{\linewidth}}}]{
                \texttt{\tiny{[P1$\langle$GM]}}
                \texttt{You are an expert in causal inference. The following question is not a typical commonsense query, but rather a meticulously designed question created by a professor specializing in causal inference, intended to assess the students' mastery of the course content.} \\
\texttt{Imagine a self{-}contained, hypothetical world with only the following conditions, and without any unmentioned factors or causal relationships: Poverty has a direct effect on water quality and cholera. Water company has a direct effect on water quality. Water quality has a direct effect on cholera. Poverty is unobserved. We know that high poverty or global water company causes clean water. high poverty or clean water causes cholera contraction. We observed the person is served by a global water company and the region has low poverty. Would the person gets cholera if polluted water instead of clean water?} \\
\texttt{Start your answer with 'Yes' or 'No', followed by additional reasoning or evidence to support your explanation.} \\
            }
        }
    }
     \\ \\

    \theutterance \stepcounter{utterance}  
    & \multicolumn{4}{p{0.6\linewidth}}{
        \cellcolor[rgb]{0.561,0.737,0.561}{
            \makecell[{{p{\linewidth}}}]{
                \texttt{\tiny{[P1$\rangle$GM]}}
                \texttt{No} \\
\\ 
\texttt{**Reasoning:**} \\
\\ 
\texttt{To determine the answer, we must analyze the causal structure and the specific counterfactual query within the constraints of the hypothetical world.} \\
\\ 
\texttt{1.  **Causal Structure Analysis:**} \\
\texttt{    *   **Parents of Water Quality (\$W\$):** Poverty (\$P\$) and Water Company (\$C\$).} \\
\texttt{    *   **Parents of Cholera (\$H\$):** Poverty (\$P\$) and Water Quality (\$W\$).} \\
\texttt{    *   **Unobserved Variable:** Poverty (\$P\$) is unobserved.} \\
\texttt{    *   **Known Causal Mechanisms:**} \\
\texttt{        *   \$P \textbackslash{}to W\$ and \$C \textbackslash{}to W\$} \\
\texttt{        *   \$P \textbackslash{}to H\$ and \$W \textbackslash{}to H\$} \\
\\ 
\texttt{2.  **Observed State:**} \\
\texttt{    *   Water Company (\$C\$) = Global (Clean water provider).} \\
\texttt{    *   Poverty (\$P\$) = Low.} \\
\texttt{    *   In this world, since \$P\$ is Low and \$C\$ is Global, the water quality (\$W\$) is Clean. Consequently, the probability of Cholera (\$H\$) is low (based on the rule "clean water causes cholera contraction" being false or low probability, while "high poverty or clean water causes cholera" implies clean water *prevents* it or is not a cause; the prompt phrasing "high poverty or clean water causes cholera contraction" is} \\
            }
        }
    }
    & & \\ \\

    \theutterance \stepcounter{utterance}  
    & & & \multicolumn{2}{p{0.3\linewidth}}{
        \cellcolor[rgb]{0.9,0.9,0.9}{
            \makecell[{{p{\linewidth}}}]{
                \texttt{\tiny{[GM$|$GM]}}
                \texttt{No} \\
            }
        }
    }
    & & \\ \\

    \theutterance \stepcounter{utterance}  
    & & & \multicolumn{2}{p{0.3\linewidth}}{
        \cellcolor[rgb]{0.9,0.9,0.9}{
            \makecell[{{p{\linewidth}}}]{
                \texttt{\tiny{[GM$|$GM]}}
                \texttt{no} \\
            }
        }
    }
    & & \\ \\

    \theutterance \stepcounter{utterance}  
    & & & \multicolumn{2}{p{0.3\linewidth}}{
        \cellcolor[rgb]{0.9,0.9,0.9}{
            \makecell[{{p{\linewidth}}}]{
                \texttt{\tiny{[GM$|$GM]}}
                \texttt{game\_result = WIN} \\
            }
        }
    }
    & & \\ \\

\end{supertabular}
}

\end{document}
